% ======================================================================
%  sec13.tex —— 第 13 节 拟凝聚层
%  底本：原书 pp.67–71（印刷页 = PDF 页）
%  原书标题：13. Quasi-coherent modules
%  体例：\section + \subsection 与原文编号一致；定理类环境紧跟 \label{类型:编号}。
% ======================================================================

\section{拟凝聚层}\label{sec:13}

% ---------------------------------------------------------------- p.67
\subsection{$\widetilde{M}$}\label{sec:13.1}

本小节中，设 $A\in\CRing$ 是交换环，$X:=\Spec(A)$ 是对应的仿射概形。考虑环化空间 $X':=(*,A)$：其底层拓扑空间是单点集，结构层由 $A\in\CRing\simeq\Shv(*,\CRing)$ 给出。设
\[
\pi\colon X\to(*,A)=X'
\]
是这样的态射：$\alpha\colon\cO_{X'}\to\pi_*\cO_X$ 对应于典范同构 $A\simeq\cO_X(X)$。注意可以作下述等同
\[
A\text{-}\mathsf{mod}\simeq\cO_{X'}\text{-}\mathsf{mod}.
\]

\begin{definition}\label{def:13.1.1}
考虑函子
\[
\pi^*\colon A\text{-}\mathsf{mod}\to\omod.
\]
对 $M\in A\text{-}\mathsf{mod}$，记
\[
\widetilde{M}:=\pi^*(M)\in\omod,
\]
并称之为与 $A$-模 $M$ 相伴的 $\cO_X$-模。
\end{definition}

\begin{yizhu}
Stacks 把同一构造写成三条等价途径：对环同态 $R\to\Gamma(X,\cO_X)$ 与 $R$-模 $M$，取 $\pi^*M$、取自由表现之后的余核、以及把 $U\mapsto\cO_X(U)\otimes_RM$ 层化，三者典范同构，结果记 $\cF_M$。该构造还与任意余极限交换，这一点在原书中要等到推论 \ref{cor:13.1.8} 才出现。
[来源：Stacks Project, Tag 01BH（Lemma 17.10.5），\url{https://stacks.math.columbia.edu/tag/01BH}（2026-11-17 访问）。]
\end{yizhu}

由推论 \ref{cor:12.3.11}，有：

\begin{corollary}\label{cor:13.1.2}
设 $A\in\CRing$，$X:=\Spec(A)$。有典范伴随
\[
\begin{array}{ccc}
A\text{-}\mathsf{mod} & \rightleftarrows & \omod\\[0.5ex]
M & \mapsto & \widetilde{M}\\
F(X) & \mapsfrom & F.
\end{array}
\]
\end{corollary}

\begin{construction}\label{con:13.1.3}
设 $A\to B$ 是 $\CRing$ 中的同态，$f\colon Y\to X$ 是对应的仿射概形之间的态射。有环化空间的典范交换方块
\[
\begin{tikzcd}[row sep=2.6em,column sep=3.4em]
Y \arrow[r] \arrow[d] & (*,B) \arrow[d]\\
X \arrow[r] & (*,A).
\end{tikzcd}
\]
由构造 \ref{con:12.3.12}，有函子的典范交换方块
\begin{equation}\tag{13.1}\label{eq:13.1}
\begin{tikzcd}[row sep=3em,column sep=3.6em]
A\text{-}\mathsf{mod} \arrow[r,"(-)(X)"] & \omod \arrow[d,"f_*"]\\
B\text{-}\mathsf{mod} \arrow[r,"(-)(Y)"'] \arrow[u,"\mathrm{res}"'] & \cO_Y\text{-}\mathsf{mod},
\end{tikzcd}
\end{equation}
以及左伴随函子的交换方块
\begin{equation}\tag{13.2}\label{eq:13.2}
\begin{tikzcd}[row sep=3em,column sep=3.6em]
A\text{-}\mathsf{mod} \arrow[r,"\widetilde{(-)}"] \arrow[d,"B\otimes_A-"'] & \omod \arrow[d,"f^*"]\\
B\text{-}\mathsf{mod} \arrow[r,"\widetilde{(-)}"'] & \cO_Y\text{-}\mathsf{mod}.
\end{tikzcd}
\end{equation}
\end{construction}

% ---------------------------------------------------------------- p.68
另一方面，用命题–定义 4.2.3 证明中的方法，可以证明下述结果：

\begin{proposition}\label{prop:13.1.4}
设 $A\in\CRing$，$X:=\Spec(A)$。对任意 $M\in A\text{-}\mathsf{mod}$，存在本质唯一的 $\cO_X$-模 $\cM$，连同 $A$-线性同构
\[
M\xrightarrow{\simeq}\cM(X),
\]
使得对任意 $f\in A$，$A$-线性映射 $M\simeq\cM(X)\to\cM(U(f))$ 诱导同构
\[
M_f\xrightarrow{\simeq}\cM(U(f)).
\]
\end{proposition}

设 $\cM$ 与 $\widetilde{M}$ 是上述 $\cO_X$-模层。按定义，$\widetilde{M}$ 是 $\pi^{\sharp}_{\PShv}(M)$ 的层化。特别地，对任意 $f\in A$，有典范的 $\cO_X(U(f))$-线性映射
\[
\cO_X(U(f))\otimes_AM\simeq(\pi^{\sharp}_{\PShv}(M))(U(f))\to\widetilde{M}(U(f)).
\]
回忆 $A_f\simeq\cO_X(U(f))$ 以及
\[
\cM(U(f))\simeq M_f\simeq A_f\otimes_AM.
\]
这给出典范的 $A_f$-线性映射
\[
\cM(U(f))\to\widetilde{M}(U(f)).
\]
由习题 1.2.7，存在唯一的态射 $\cM\to\widetilde{M}$（在 $\omod$ 中），它在每个 $U(f)$ 上的取值由上述映射给出。

\begin{proposition}\label{prop:13.1.5}
设 $A\in\CRing$，$X:=\Spec(A)$。对 $M\in A\text{-}\mathsf{mod}$，设 $\widetilde{M}$ 与 $\cM$ 是上述 $\cO_X$-模。则典范态射 $\cM\to\widetilde{M}$ 是同构。特别地，对任意 $f\in A$，有
\[
M_f\xrightarrow{\simeq}\widetilde{M}(U(f)).
\]
\end{proposition}

\begin{proof}
设 $\cN\in\omod$ 是测试对象，只需说明映射
\begin{equation}\tag{13.3}\label{eq:13.3}
\Hom_{\omod}(\widetilde{M},\cN)\to\Hom_{\omod}(\cM,\cN)
\end{equation}
是双射。按层化的定义，有
\begin{equation}\tag{13.4}\label{eq:13.4}
\Hom_{\omod}(\widetilde{M},\cN)\xrightarrow{\simeq}\Hom_{\omodP}(\pi^{\sharp}_{\PShv}(M),\cN).
\end{equation}
注意 $\pi^{\sharp}_{\PShv}(M)$ 与 $\cM$ 在 $U(f)$ 上的取值都与 $M_f$ 有典范等同。由习题 1.2.7，得到等同
\begin{equation}\tag{13.5}\label{eq:13.5}
\Hom_{\omod}(\cM,\cN)\simeq\Hom_{\omodP}(\pi^{\sharp}_{\PShv}(M),\cN).
\end{equation}
容易验证 \eqref{eq:13.4} 等于复合 \eqref{eq:13.5}$\circ$\eqref{eq:13.3}。于是 \eqref{eq:13.3} 也是双射。
\end{proof}

\begin{yizhu}
命题的另一说法是泛性质 $\Hom_{\cO_X}(\widetilde{M},\cF)\simeq\Hom_A(M,\Gamma(X,\cF))$：$\widetilde{M}$ 由 $M$ 与「取整体截面」这一右伴随共同决定。等式 $M_f\simeq\widetilde{M}(U(f))$ 正是该泛性质在主开集上的显式形式，也是把仿射概形上的 $\cO_X$-模等同于 $A$-模的全部内容。
[来源：Stacks Project, Tag 01I7（Lemma 26.7.1），\url{https://stacks.math.columbia.edu/tag/01I7}（2026-11-17 访问）。]
\end{yizhu}

\begin{corollary}\label{cor:13.1.6}
设 $A\in\CRing$，$X:=\Spec(A)$。函子
\[
\widetilde{(-)}\colon A\text{-}\mathsf{mod}\to\omod
\]
全忠实。
\end{corollary}

\begin{proof}
该函子有右伴随 $\cF\mapsto\cF(X)$，且单位自然变换 $M\to\widetilde{M}(X)$ 可逆（命题 \ref{prop:13.1.5}）。由此形式地推出左伴随全忠实。
\end{proof}

% ---------------------------------------------------------------- p.69
\begin{warning}\label{warn:13.1.7}
函子 $\widetilde{(-)}\colon A\text{-}\mathsf{mod}\to\omod$ 不本质满射。例如，习题 11.4.6 中 $\cO_X$ 的理想层 $\cI$ 不同构于任何 $\widetilde{M}$（$M\in k[t]\text{-}\mathsf{mod}$），因为 $\cI(X)\simeq 0$。
\end{warning}

\begin{yizhu}
从 stack 的观点看，$\cO_X$-模层构成 $X$ 上的一个 stack，拟凝聚层是其中「局部自由层之间态射的余核」那一片。nLab 的条目把这一点列为拟凝聚层的等价刻画，并指出它也可说成从 $\CRing$ 到 $\Mod$ 的、与拉回相容的截面。$\widetilde{(-)}$ 不满射，正说明这类截面比任意模层要少。
[来源：nLab "quasicoherent sheaf"，\url{https://ncatlab.org/nlab/show/quasicoherent+sheaf}（Last revised 2026-06-15；2026-11-17 访问）。]
\end{yizhu}

\begin{corollary}\label{cor:13.1.8}
设 $A\in\CRing$，$X:=\Spec(A)$。函子
\[
\widetilde{(-)}\colon A\text{-}\mathsf{mod}\to\omod
\]
与有限极限及小余极限交换。
\end{corollary}

\begin{proof}
由命题 \ref{prop:13.1.5} 推出，并用 $M\mapsto M_f$ 与有限极限、小余极限交换。
\end{proof}

\begin{corollary}\label{cor:13.1.9}
设 $\phi\colon A\to B$ 是 $\CRing$ 中的同态，$f\colon Y\to X$ 是对应的仿射概形之间的态射。交换方块 \eqref{eq:13.1} 诱导自然同构 $\widetilde{\mathrm{res}(-)}\simeq f_*\widetilde{(-)}$，即有交换方块
\begin{equation}\tag{13.6}\label{eq:13.6}
\begin{tikzcd}[row sep=3em,column sep=3.6em]
A\text{-}\mathsf{mod} \arrow[r,"\widetilde{(-)}"] & \omod \arrow[d,"f_*"]\\
B\text{-}\mathsf{mod} \arrow[r,"\widetilde{(-)}"'] \arrow[u,"\mathrm{res}"'] & \cO_Y\text{-}\mathsf{mod}.
\end{tikzcd}
\end{equation}
\end{corollary}

\begin{proof}
对 $M\in A\text{-}\mathsf{mod}$ 与 $f\in A$，命题 \ref{prop:13.1.5} 给出
\[
\widetilde{\mathrm{res}(M)}(U(f))\simeq\mathrm{res}(M)_f\simeq M_{\phi(f)}\simeq\widetilde{M}(U(\phi(f)))\simeq f_*\widetilde{M}(U(f)).
\]
容易验证，这一复合正是 Beck--Chevalley 自然变换 $\widetilde{\mathrm{res}(-)}\to f_*\widetilde{(-)}$ 在 $U(f)$ 上的取值。由此得到所要结论。
\end{proof}

% ---------------------------------------------------------------- p.69–70
\subsection{拟凝聚层的定义}\label{sec:13.2}

\begin{definition}\label{def:13.2.1}
设 $X$ 是概形。称 $\cO_X$-模 $\cF$ \emph{拟凝聚}，若对任意开浸入 $j\colon U=\Spec(A)\to X$，$\cO_U$-模 $\cF|_U:=j^*\cF$ 同构于某个 $\widetilde{M}$（$M\in A\text{-}\mathsf{mod}$）。

记
\[
\mathsf{QCoh}(X)=\omod_{\mathrm{qcoh}}\subseteq\omod
\]
为拟凝聚 $\cO_X$-模构成的全子范畴。
\end{definition}

\begin{yizhu}
Stacks 的定义不依赖仿射开集：$\cF$ 拟凝聚，指每点有邻域 $U$，使 $\cF|_U$ 是某个态射 $\bigoplus_J\cO_U\to\bigoplus_I\cO_U$ 的余核。该定义对任意环化空间都可用，且一般环化空间上拟凝聚层范畴未必是阿贝尔范畴。概形上两种定义一致：仿射开集上 $\widetilde{M}$ 恰是自由模表现余核的层化。
[来源：Stacks Project, Tag 01BD（Definition 17.10.1），\url{https://stacks.math.columbia.edu/tag/01BD}（2026-11-17 访问）。]
\end{yizhu}

\begin{example}\label{ex:13.2.2}
结构层 $\cO_X$ 作为自身的模是拟凝聚的。
\end{example}

% ---------------------------------------------------------------- p.69–70
\subsection{仿射情形}\label{sec:13.3}

\begin{proposition}\label{prop:13.3.1}
设 $A\in\CRing$，$X:=\Spec(A)$。下述函子有定义，且互为逆函子
\[
\begin{array}{ccc}
A\text{-}\mathsf{mod} & \rightleftarrows & \mathsf{QCoh}(X)\\[0.5ex]
M & \mapsto & \widetilde{M}\\
\cF(X) & \mapsfrom & \cF.
\end{array}
\]
\end{proposition}

\begin{proof}
先证有定义：只需说明 $\widetilde{M}$ 拟凝聚（$M\in A\text{-}\mathsf{mod}$）。设 $U=\Spec(B)\to X$ 是由同态 $A\to B$ 给出的开浸入。由构造 \ref{con:13.1.3}，有 $\widetilde{M}|_U\simeq\widetilde{B\otimes_AM}$。按定义，这说明 $\widetilde{M}$ 拟凝聚。

由推论 \ref{cor:13.1.2}，只需说明函子 $A\text{-}\mathsf{mod}\to\mathsf{QCoh}(X)$、$M\mapsto\widetilde{M}$ 是等价。由推论 \ref{cor:13.1.6}，只需说明它本质满射。而这一点由定义即得。
\end{proof}

\begin{yizhu}
这条等价连模的同态也一并算进去了：$M$ 有限表现时，$\uHom_{\cO_X}(\widetilde{M},\widetilde{N})\simeq\widetilde{\Hom_R(M,N)}$。左端是层，右端由模层化而来，故该式说明「取整体截面」不仅给出对象之间的一一对应，还与内 Hom 相容。这也解释了为何仿射概形上的同调计算可以整段搬回交换代数。
[来源：Stacks Project, Tag 01I6（Section 26.7, Quasi-coherent sheaves on affines），\url{https://stacks.math.columbia.edu/tag/01I6}（2026-11-17 访问）。]
\end{yizhu}

结合命题 \ref{prop:13.3.1} 与推论 \ref{cor:13.1.8}，得到：

\begin{corollary}\label{cor:13.3.2}
设 $X$ 是仿射概形。拟凝聚 $\cO_X$-模的有限极限与小余极限仍拟凝聚。
\end{corollary}

结合命题 \ref{prop:13.3.1} 与 \eqref{eq:13.2}，得到：

\begin{corollary}\label{cor:13.3.3}
设 $f\colon X\to Y$ 是仿射概形之间的态射。函子 $f^*$ 把拟凝聚 $\cO_Y$-模送到拟凝聚 $\cO_X$-模。
\end{corollary}

结合命题 \ref{prop:13.3.1} 与推论 \ref{cor:13.1.9}，得到：

\begin{corollary}\label{cor:13.3.4}
设 $f\colon X\to Y$ 是仿射概形之间的态射。函子 $f_*$ 把拟凝聚 $\cO_X$-模送到拟凝聚 $\cO_Y$-模。
\end{corollary}

% ---------------------------------------------------------------- p.70
\subsection{一般情形}\label{sec:13.4}

\begin{theorem}\label{thm:13.4.1}
拟凝聚性是局部条件。换言之，对概形 $X$ 与开覆盖 $X=\bigcup_{i\in I}X_i$，对象 $\cF\in\omod$ 拟凝聚，当且仅当每个 $\cF|_{X_i}$（$i\in I$）拟凝聚。
\end{theorem}

\begin{proof}
「仅当」显然。

对「若」，设 $\cF$ 是 $\cO_X$-模，且每个 $\cF|_{X_i}$ 拟凝聚。只需说明对任意仿射开子概形 $U\subseteq X$，$\cF|_U$ 拟凝聚。取 $U$ 的仿射开子概形覆盖 $U=\bigcup_{j\in J}U_j$，使每个 $U_j$ 都含于某个 $X_i$。由假设，每个 $\cF|_{U_j}\simeq(\cF|_{X_i})|_{U_j}$ 拟凝聚。于是可以把 $(X,\{X_i\},\cF)$ 换成 $(U,\{U_j\},\cF|_U)$，化归到 $X$ 与各 $X_i$ 皆仿射的情形。

$X$ 仿射时，其底层拓扑空间拟紧，故可设 $I$ 是有限集。记 $X_{ij}:=X_i\cap X_j$，并设 $f_i\colon X_i\to X$、$f_{ij}\colon X_{ij}\to X$ 是开浸入。$\cF$ 是层，故由习题 1.3.4 有同构
\[
\cF\simeq\ker\Bigl(\prod_{i\in I}f_{i,*}(\cF|_{X_i})\to\prod_{(i,j)\in I^2}f_{ij,*}(\cF|_{X_{ij}})\Bigr).
\]
按假设，每个 $\cF|_{X_i}$ 拟凝聚，故每个 $\cF|_{X_{ij}}$ 也拟凝聚。由推论 \ref{cor:13.3.4}，$f_{i,*}(\cF|_{X_i})$ 与 $f_{ij,*}(\cF|_{X_{ij}})$ 都拟凝聚。由推论 \ref{cor:13.3.2}，$\cF$ 拟凝聚。
\end{proof}

\begin{corollary}\label{cor:13.4.2}
设 $f\colon X\to Y$ 是概形之间的态射。函子 $f^*$ 把拟凝聚 $\cO_Y$-模送到拟凝聚 $\cO_X$-模。
\end{corollary}

% ---------------------------------------------------------------- p.71
\begin{proof}
用定理 \ref{thm:13.4.1}，把 $X\to Y$ 换成它在仿射开子概形 $U\subseteq Y$ 上的基变换 $X\times_YU\to U$，可设 $Y$ 仿射。再用定理 \ref{thm:13.4.1}，把 $X$ 换成它的仿射开子概形 $V$，可设 $X$ 仿射。此时结论由推论 \ref{cor:13.3.3} 给出。
\end{proof}

\begin{corollary}\label{cor:13.4.3}
设 $X$ 是概形。拟凝聚 $\cO_X$-模的有限极限与小余极限仍拟凝聚。
\end{corollary}

\begin{proof}
由推论 \ref{cor:13.3.2} 与下述事实推出：对任意开浸入 $j$，函子 $j^{-1}$ 与有限极限、小余极限交换。
\end{proof}

\begin{lemma}\label{lem:13.4.4}
设 $f\colon X\to Y$ 是闭浸入。函子 $f_*$ 把拟凝聚 $\cO_X$-模送到拟凝聚 $\cO_Y$-模。
\end{lemma}

\begin{proof}
问题在 $Y$ 上局部，故可设 $Y$ 仿射，从而拟紧。$f$ 是闭浸入，故 $X$ 拟紧。取 $X$ 的有限仿射开子概形覆盖 $X=\bigcup_{i\in I}X_i$。有
\[
X_{ij}:=X_i\cap X_j\simeq X_i\times_XX_j\xrightarrow{\simeq}X_i\times_YX_j,
\]
最后一个同构是因为 $f\colon X\to Y$ 是单态射（命题 11.3.7）。特别地，$X_{ij}$ 仿射。设 $g_i\colon X_i\to X$、$g_{ij}\colon X_{ij}\to X$ 是开浸入。

设 $\cF\in\mathsf{QCoh}(X)$。只需说明 $f_*\cF$ 拟凝聚。与定理 \ref{thm:13.4.1} 的证明相同，有
\[
\cF\simeq\ker\Bigl(\prod_{i\in I}g_{i,*}(\cF|_{X_i})\to\prod_{(i,j)\in I^2}g_{ij,*}(\cF|_{X_{ij}})\Bigr).
\]
$f_*$ 是右伴随，故与极限交换。这给出
\[
f_*\cF\simeq\ker\Bigl(\prod_{i\in I}(f\circ g_i)_*(\cF|_{X_i})\to\prod_{(i,j)\in I^2}(f\circ g_{ij})_*(\cF|_{X_{ij}})\Bigr).
\]
由推论 \ref{cor:13.4.3} 与推论 \ref{cor:13.3.4} 即得结论。
\end{proof}

\begin{remark}\label{rem:13.4.5}
后续课程中会证明：若 $f$（相对地）拟紧且拟分离，则 $f_*$ 保持拟凝聚模。
\end{remark}

\begin{yizhu}
注记 \ref{rem:13.4.5} 预告的结论在 Stacks 中有完整表述：$f$ 拟紧且拟分离时，任意拟凝聚 $\cF$ 的高阶正像 $R^pf_*\cF$ 在目标上仍拟凝聚。证明沿目标的仿射开覆盖化归，用的正是本节这套局部化手法，只是把核换成了 Čech 复形。
[来源：Stacks Project, Tag 01XJ（Lemma 30.4.5），\url{https://stacks.math.columbia.edu/tag/01XJ}（2026-11-17 访问）。]
\end{yizhu}

\begin{warning}\label{warn:13.4.6}
对一般的态射 $f\colon X\to Y$，$f_*$ 不保持拟凝聚模。
\end{warning}

\begin{exercise}\label{exc:13.4.7}
设 $k$ 是域，$X:=\A^1_k$。
\begin{enumerate}[label=(\arabic*),leftmargin=2em,itemsep=0.25ex]
  \item 证明拟凝聚 $\cO_X$-模的无穷积可能不拟凝聚。
  \item 设 $I$ 是无限集，考虑显然的态射
  \[
  (\id)_{i\in I}\colon\bigsqcup_{i\in I}\A^1_k\to\A^1_k.
  \]
  证明沿该态射的前推不保持拟凝聚模。
\end{enumerate}
\end{exercise}

\begin{exercise}\label{exc:13.4.8}
设 $X$ 是概形，$\cM,\cN\in\mathsf{QCoh}(X)$。证明 $\cM\otimes_{\cO_X}\cN\in\mathsf{QCoh}(X)$。
\end{exercise}
